% Imporditud: tools/import_eksperimendid.py
% Allikas: Eesti füüsikaolümpiaadi eksperimendi ülesannete
%   kogu (Rael Kalda, 2023), ülesanne Ü75, lahendus L75.
\setAuthor{EFO žürii}
\setRound{lõppvoor}
\setYear{1996}
\setNumber{G E1}
\setDifficulty{3}
\setTopic{vedelike mehaanika}
\setVahendid{kaks keha, mõõtejoonlaud, pliiats, nihik}
\setPunktid{10}
\setAllikas{Eksperimendi kogu Ü75, lahendus lk 67}
\setLahendusseis{toores}

\prob{Kahe keha ruumalad}
Määrata kahe samast materjalist keha ruumala. Esitada katse idee. Kirjeldada tegevuse etappe. Hinnata mõõteviga.

\hint

\solu
Korrapärase kujuga keha ruumala määramine (20\%) (30\%) Idee korrapäratu kujuga keha ruumala määramiseks: $m_1$ = $V_1$

m$^2$ $V_2$ Masside suhte määramine (20\%) Korrapäratu keha ruumala leidmine (10\%) Mõõtmistäpsuse analüüs (20\%)
\probend
